%% kapa: exam=1 round=31 year=2020 period=1 form=A subject=civil_law subject_ko=민법 seq=29 number=29
%% source: Q-Net 감정평가사 기출문제 (cst00309) artlSeq=5207219 file=extracted/1008/5207219/A형/A형/2020년도 제31회 감정평가사 1차 시험_1교시_A형.pdf
%% conversion: kordoc 4.24.51 -> tools/gen1.py (auto); review=proofread-sonnet; flags=none
\begin{question}{29}
소유권에 관한 설명으로 옳지 않은 것은? (다툼이 있으면 판례에 따름)
\begin{choices}
  \item 매도인은 매매계약의 이행으로 토지를 인도받았으나 소유권이전등기를 하지 않고 점유ㆍ사용하는 매수인에게 부당이득의 반환을 청구할 수 있다.
  \item 토지의 경계는 지적공부에 의하여 확정된다.
  \item 토지가 포락되어 사회통념상 원상복구가 어려워 토지로서의 효용을 상실한 때에는 그 토지의 소유권이 소멸한다.
  \item 도급계약에서 신축집합건물의 소유권을 수인의 도급인에게 귀속할 것을 약정한 경우, 그 건물의 각 전유부분의 소유관계는 공동도급인의 약정에 의한다.
  \item 소유권에 기한 방해배제청구권에서 ‘방해’는 현재 지속되고 있는 침해를 의미한다.
\end{choices}
\end{question}
\answer{1}
